\section*{Chapter 5: Conclusion and Further Work}
\phantomsection \label{chap:conclusion} \setcounter{section}{5} \setcounter{subsection}{0}
\addcontentsline{toc}{section}{\textit{Chapter 5: Conclusion and Further Work}}

\subsection{Conclusion}

This project set out to identify which parts of an MCP-style control 
plane genuinely benefit from kernel-level enforcement. 
The answer is 
a small and well-defined set: semantic identity, requester binding, 
approval validity, and invocation completion. 
These four invariants 
are structurally difficult to hold when enforcement is confined to a 
user-space gateway.

\texttt{linux-mcp} moves only those invariants into a Linux kernel 
module, while leaving semantic parsing and tool execution in user 
space. 
The gateway remains fully trusted, but the state it establishes 
now resides in the kernel, beyond the reach of unprivileged processes 
and durable across daemon failure.

In the tested prototype setting, this partitioning blocked the full 
set of 110 targeted control-plane attack cases in our functional 
evaluation, where both user-space baselines permitted meaningful 
bypass rates.
Latency remained statistically indistinguishable from 
the user-space baseline at small and medium payloads, was measurably 
lower at 1\,MB, and throughput stayed within approximately 10\% of 
the user-space peak. \texttt{linux-mcp} is therefore both a working 
prototype and a design claim: the enforcement gains are bounded in 
cost and precise in scope.

\subsection{Reflection}

The most important question clarified by this project was not simply ``where should code run?'', but ``which invariants must remain trustworthy even when a gateway crashes, is misconfigured, or behaves incorrectly?''
This reframing shaped the entire design: rather than asking how to make the gateway more secure, the project asked which properties the gateway should not be solely responsible for enforcing.

From a technical standpoint, the project required working across layers that are often studied separately: implementing a Linux kernel module, using Generic Netlink as a structured kernel--user-space control interface, exporting live state through \texttt{sysfs}, and coordinating these pieces with a manifest-aware daemon and UDS tool services.
This process deepened understanding of interface ownership, failure handling, recovery logic, and the practical cost of crossing privilege boundaries.
A central lesson of systems work was reinforced: the harder problem is often not adding mechanisms, but deciding what each layer should and should not be responsible for.

A second lesson concerned research discipline.
Early iterations of this project were formulated around broad claims such as ``putting security in the kernel'' or ``making agents safer''.
Such claims lack sufficient specificity to be empirically testable.
The project became considerably stronger once the problem was narrowed to four concrete control-plane properties (semantic identity binding, approval validity, hash consistency, and durable observability), at which point the threat model, implementation choices, and evaluation criteria aligned naturally.
A systems argument becomes credible only when the claimed benefit is specific enough to be tested directly.

The project also reinforced the value of restraint in privileged-system design.
The temptation to move more logic into the kernel was real: every additional check the kernel performs is one the gateway cannot bypass.
But manifest parsing, schema validation, and tool execution all belong in user space, where iteration is fast and debugging does not require a VM reboot.
A privileged component that does too much is harder to reason about and harder to trust---the small kernel role here was not a limitation but a feature.

One practical consequence of this project is an appreciation of how much the framing of a security claim matters.
\texttt{linux-mcp} improves a specific part of the control plane; it does not address unsafe tool semantics, poor policy choices, or compromise outside the assumed threat model.
Getting the scope right was not just a methodological exercise: a system that claims to solve the full AI-agent security problem is harder to evaluate, harder to trust, and ultimately harder to deploy than one that solves a precise sub-problem well.
The hardest part was resisting the temptation to widen the scope as the implementation took shape.

The broader lesson is simply that better security is not always achieved by adding more machinery.
Moving a small set of authority-bearing state transitions into a stronger enforcement domain turned out to be more tractable, and arguably more useful, than attempting a fuller redesign.

\subsection{Further Work}
\label{sec:further-work}

Several directions follow naturally from the current prototype.

First, the current kernel arbitration logic is intentionally minimal.
It verifies tool identity, agent identity, semantic-hash consistency, and approval-ticket validity, but it does not implement a richer policy framework.
A natural next step is to support more expressive admission policies, such as per-tool permissions, role-scoped access control, rate limits (e.g., $N$ requests per minute per agent), or execution budgets (e.g., total time per session).
The design principle should be maintained: the kernel verifies policy membership but does not compute policy decisions, which remain in user space, preserving the narrow-kernel invariant.

Second, the approval mechanism can be developed into a more complete governance path. 
The present design supports defer-and-decide control for sensitive operations, but real deployments may require approval delegation, bounded approval scope, revocation, expiry policies, and richer audit trails linking each decision to a concrete execution context. 
Strengthening this path would make \texttt{linux-mcp} more suitable for environments with accountability or compliance requirements.

Third, recovery behaviour can be improved further. The current prototype already preserves kernel-held registry and approval state across daemon failure, but reconciliation is deliberately conservative and session-level state is rebuilt after restart. 
A more mature design could provide more systematic restart handling, richer recovery metadata, and clearer support for interrupted or partially completed requests.
This becomes particularly important for long-running agent workflows, where failure and restart are routine rather than exceptional.

Fourth, the evaluation can be extended substantially.
Future work should include bare-metal evaluation to eliminate hypervisor timing noise, concurrency levels well beyond the current maximum of 100 concurrent requests (e.g., 500 or 1000), tool registries an order of magnitude larger than those used here, and heterogeneous tool workloads with variable or network-dependent execution times.
$p_{95}$ latency is characterised in Chapter~\ref{chap:evaluation}; $p_{99}$ under sustained concurrency remains unmeasured and is a natural next step.
Comparison against additional hardened user-space designs would also help sharpen the argument by separating the benefits of kernel placement from those achievable through improved user-space engineering alone.

Finally, this work opens a broader research question: which other agent-oriented abstractions deserve operating-system support?
\texttt{linux-mcp} shows that a small but consequential part of the MCP control plane benefits from OS-level support without requiring AI agents to become first-class OS citizens or requiring a new OS substrate.
This project focuses on admission control for tool invocation, but similar arguments may extend to execution budgets, provenance across tool chains, durable approval logs, or agent-scoped resource accounting.
Exploring these directions would help clarify whether support for AI agents should be understood as entirely new operating-system substrates, or as conventional systems extended with a small number of carefully scoped, agent-aware mechanisms.
In that sense, the contribution of \texttt{linux-mcp} is best read as a precise first step rather than a complete answer.
